% Nature COVID-19 RBC Microangiopathy 2025 教學講義
% 缺血性內皮細胞 necroptosis 誘導紅血球溶血與 COVID-19 血管病變
\documentclass[12pt,a4paper]{article}
% Drake_preamble.tex - LaTeX Preamble Template
% 作者: 謝慕揚, MD, PhD, FESC
% 
% 使用說明:
% 1. 表格請使用 longtable，不要有垂直分隔線 |
% 2. 表格內換行使用 \newline，不要使用 \\
% 3. 藥物名稱、檢驗、檢查請維持英文
% 4. Reference 格式請使用 \href 加上驗證過的 DOI

% Including essential packages for Chinese typesetting
\usepackage{xeCJK}
\usepackage{fontspec}
% Configuring Chinese font with Noto Serif CJK TC, fallback to SimSun
\IfFontExistsTF{Noto Serif CJK TC}{
  \setCJKmainfont{Noto Serif CJK TC}
  \setCJKsansfont{Noto Serif CJK TC}
  \setCJKmonofont{Noto Serif CJK TC}
}{\setCJKmainfont{SimSun}
  \setCJKsansfont{SimSun}
  \setCJKmonofont{SimSun}
}

% Including other necessary packages
\usepackage{geometry}
\usepackage{titlesec}
\usepackage{enumitem}
\usepackage{xcolor}
\usepackage{colortbl}
\usepackage{tcolorbox}
\usepackage{booktabs}
\usepackage{longtable}
\usepackage{array}
\usepackage{graphicx}
\usepackage{tikz}
\usepackage{fancyhdr}
\usepackage{multicol}
\usepackage{datetime2}
\usepackage{hyperref}
\usepackage{mdframed}
\usepackage{amsmath}
\usepackage{amssymb}
\usepackage{multirow}

\hypersetup{
    colorlinks=true,
    linkcolor=emergency_blue,
    citecolor=emergency_blue,
    urlcolor=emergency_blue,
    pdfborder={0 0 0}
}

% Defining page geometry
\geometry{
  top=2.5cm,
  bottom=2.5cm,
  left=2cm,
  right=2cm,
  headheight=25pt
}

% Adjusting topmargin to compensate for increased headheight
\addtolength{\topmargin}{-10pt}

% Defining colors
\definecolor{emergency_red}{RGB}{186,24,27}
\definecolor{emergency_blue}{RGB}{0,114,188}
\definecolor{emergency_gray}{RGB}{120,120,120}
\definecolor{highlight_yellow}{RGB}{255,250,205}

% Setting title formats
\titleformat{\section}[block]{\Large\bfseries\color{emergency_red}}{\thesection}{10pt}{}
\titleformat{\subsection}[block]{\large\bfseries\color{emergency_blue}}{\thesubsection}{8pt}{}
\titleformat{\subsubsection}[block]{\normalsize\bfseries\color{emergency_gray}}{\thesubsubsection}{6pt}{}

% Defining custom environment for important notes
\newmdenv[
    linecolor=emergency_red,
    backgroundcolor=highlight_yellow,
    linewidth=2pt,
    topline=true,
    bottomline=true,
    leftline=true,
    rightline=true,
    innertopmargin=10pt,
    innerbottommargin=10pt,
    innerleftmargin=10pt,
    innerrightmargin=10pt
]{importantbox}

% Defining note command with tcolorbox
\newcommand{\note}[1]{
    \par
    \noindent
    \begin{tcolorbox}[
        colback=emergency_blue!5!white,
        colframe=emergency_blue,
        title=\textbf{重點提醒},
        fonttitle=\bfseries,
        boxrule=0.5pt,
        arc=3pt,
        left=6pt,
        right=6pt,
        top=6pt,
        bottom=6pt
    ]
    #1
    \end{tcolorbox}
}

% Key findings box
\newcommand{\keyfinding}[1]{
    \par
    \noindent
    \begin{tcolorbox}[
        colback=yellow!10!white,
        colframe=orange!75!black,
        title=\textbf{核心發現摘要},
        fonttitle=\bfseries,
        boxrule=0.5pt,
        arc=3pt
    ]
    #1
    \end{tcolorbox}
}

% Defining custom date format for ISO 8601
\DTMnewdatestyle{iso}{
  \renewcommand{\DTMdisplaydate}[4]{\number##1-\number##2-\number##3}
  \renewcommand{\DTMDisplaydate}[4]{\number##1-\number##2-\number##3}
}
\DTMsetdatestyle{iso}
\newcommand{\mydate}{\DTMtoday}


% Custom header for this document
\pagestyle{fancy}
\fancyhf{}
\fancyhead[L]{\textcolor{emergency_red}{Nature 教學講義}}
\fancyhead[R]{\textcolor{emergency_blue}{COVID-19 RBC Microangiopathy}}
\fancyfoot[C]{\textcolor{emergency_gray}{\thepage}}
\renewcommand{\headrulewidth}{1pt}
\renewcommand{\headrule}{\hbox to\headwidth{\color{emergency_red}\leaders\hrule height \headrulewidth\hfill}}

\begin{document}

%============================================================
% Title Page
%============================================================
\begin{center}
{\LARGE\bfseries\textcolor{emergency_red}{Nature}}\\[0.5cm]
{\Large\textbf{Ischaemic Endothelial Necroptosis Induces Haemolysis and COVID-19 Angiopathy}}\\[0.3cm]
{\large 缺血性內皮細胞壞死性凋亡誘導紅血球溶血與 COVID-19 血管病變}\\[0.5cm]
{\normalsize \href{https://doi.org/10.1038/s41586-025-09076-x}{Nature 2025. DOI: 10.1038/s41586-025-09076-x}}\\[0.3cm]
{\normalsize 整理：謝慕揚 MD, PhD, FESC \quad 日期：\mydate}
\end{center}

\vspace{0.5cm}

%============================================================
\section{研究背景}
%============================================================

\subsection{COVID-19 微血管病變的重要性}

COVID-19 的 microangiopathy（微血管病變）是 SARS-CoV-2 感染的重要併發症：
\begin{itemize}
    \item 造成急性及慢性併發症（long COVID syndrome）
    \item 毛細血管功能下降 60--90\%
    \item 傳統認為是由 \textbf{fibrin}（纖維蛋白）和 \textbf{platelet}（血小板）微血栓造成
    \item 但 anticoagulant therapy（抗凝血治療）效果有限
\end{itemize}

\subsection{傳統觀點的問題}

過去認為 COVID-19 微血管阻塞主要是由於：
\begin{itemize}
    \item Endotheliopathy（內皮病變）
    \item 血管收縮
    \item Microthrombus formation（微血栓形成）
    \item COVID-19-associated coagulopathy
\end{itemize}

\note{本研究挑戰了傳統觀點，發現 COVID-19 微血管阻塞的主要原因並非 fibrin 和 platelet 沉積，而是一種全新的 \textbf{RBC haemolysis}（紅血球溶血）機制。}

%============================================================
\section{核心發現摘要}
%============================================================

\keyfinding{
\begin{enumerate}
    \item \textbf{廣泛的內皮細胞死亡}：COVID-19 器官微血管中存在大量 EC death
    \item \textbf{非傳統機制}：EC death 與 fibrin 或 platelet 沉積無關，而是與 RBC haemolysis 相關
    \item \textbf{缺血再灌注損傷}：此 RBC microangiopathy 與 ischaemia-reperfusion injury 相關
    \item \textbf{跨疾病相關性}：同樣機制見於 MI、腸道缺血、stroke、septic shock、cardiogenic shock
    \item \textbf{分子機制}：Ischaemia 誘導 \textbf{MLKL-dependent EC necroptosis} 和 \textbf{complement-dependent RBC haemolysis}
    \item \textbf{新型止血機制}：溶血的 RBC membranes 形成 ``endo-seal''，獨立於 platelet 和 fibrin
    \item \textbf{治療標靶}：EC Mlkl 基因剔除可減少 RBC haemolysis 和缺血性器官損傷
\end{enumerate}
}

%============================================================
\section{研究方法}
%============================================================

\subsection{多模態影像分析}

研究團隊開發了定量、多重、多模態影像方法：
\begin{itemize}
    \item 分析超過 1,000 條血管（動脈、靜脈、毛細血管）
    \item 來自 COVID-19 死亡患者的肺、心、腎、肝
    \item 使用 immunofluorescence (IF)、H\&E、scanning electron microscopy (SEM)
    \item Correlative light and electron microscopy (CLEM) 分析
\end{itemize}

\subsection{動物模型}

\begin{itemize}
    \item K18-hACE2 小鼠 + SARS-CoV-2 Delta 株感染
    \item LAD artery ligation（心臟缺血模型）
    \item Superior mesenteric artery ligation（腸道 IR 模型）
    \item 缺血性 stroke 模型
    \item Heart transplant 模型
\end{itemize}

\subsection{基因剔除小鼠}

\begin{longtable}{p{4cm}p{10cm}}
\toprule
\textbf{小鼠品系} & \textbf{用途} \\
\midrule
\endhead
Mlkl$^{-/-}$ & 全身性 MLKL 缺失，確認 necroptosis 在 RBC haemolysis 的角色 \\
Mlkl$^{EC-/-}$ & 內皮細胞特異性 MLKL 缺失，確認 EC necroptosis 的角色 \\
Mlkl$^{RBC-/-}$ & 紅血球特異性 MLKL 缺失，排除 RBC necroptosis 的角色 \\
C9$^{-/-}$ & Complement C9 缺失，確認 membrane attack complex 在 RBC haemolysis 的角色 \\
\bottomrule
\end{longtable}

%============================================================
\section{詳細結果}
%============================================================

\subsection{COVID-19 器官中的內皮細胞死亡}

\subsubsection{內皮細胞損傷的證據}

在 COVID-19 死亡患者器官中發現：
\begin{itemize}
    \item 廣泛的 VE-cadherin 流失
    \item PECAM-1 表面標記消失
    \item Glycocalyx（糖萼）破壞
    \item EC 核型態改變：pyknosis（核濃縮）、karyorrhexis（核碎裂）、karyolysis（核溶解）
    \item EC detachment（10--50\% 血管）
\end{itemize}

\subsubsection{EC 死亡與組織損傷的相關性}

\begin{importantbox}
\textbf{關鍵發現}：EC death 的程度與 tissue injury 的嚴重度呈正相關。輕度組織損傷時 EC death 較少，重度組織損傷時 EC death 明顯增加（p $<$ 0.0001）。
\end{importantbox}

\subsection{RBC Haemolysis 在 COVID-19 血管中}

\subsubsection{RBC 衍生的無定形物質}

在 EC death 部位發現廣泛沉積的 amorphous, proteinaceous acellular material：
\begin{itemize}
    \item 心臟微血管：30\%
    \item 腎臟微血管：29\%
    \item 肝臟微血管：27\%
    \item 肺臟微血管：5\%（較少）
\end{itemize}

\subsubsection{CD235（Glycophorin A）染色特徵}

\begin{longtable}{p{4cm}p{5cm}p{5cm}}
\toprule
\textbf{特徵} & \textbf{完整 RBC (CD235$^{low}$)} & \textbf{溶血 RBC (CD235$^{high}$)} \\
\midrule
\endhead
CD235 螢光強度 & 低 & 高（約 5 倍增加） \\
形態 & 正常雙凹圓盤狀 & 不規則、碎片化 \\
Haemoglobin & 保留 & 流失（Hb$^{-}$） \\
沉積位置 & 血管腔內 & 血管壁及 RBC 之間 \\
\bottomrule
\end{longtable}

\subsubsection{非 DIC 相關}

\begin{itemize}
    \item 典型的 disseminated intravascular coagulation (DIC) 組織學特徵（如腎臟毛細血管 fibrin）在 COVID-19 屍檢組織中罕見
    \item 患者無 DIC 的血液學改變
    \item RBC haemolysis 與 EC death 有強烈空間相關性
\end{itemize}

\note{這表明 COVID-19 microangiopathy 與 COVID-19-associated coagulopathy 無關，而是與發生在 EC death 部位的 RBC haemolytic process 相關。}

\subsection{器官缺血誘導 RBC Haemolysis}

\subsubsection{非 COVID-19 缺血性器官的發現}

研究團隊檢查了非 COVID-19 患者的缺血性器官：

\begin{longtable}{p{4cm}p{10cm}}
\toprule
\textbf{器官/疾病} & \textbf{發現} \\
\midrule
\endhead
Acute myocardial infarction (AMI) & 廣泛的 RBC haemolysis 和 EC death，RBC membrane 沉積於血管壁和 RBC 之間，造成微血管阻塞 \\
Ischaemic gut & 同樣模式的 RBC haemolysis 和 EC death \\
Ischaemic stroke & 同樣模式 \\
Cardiogenic shock & 多器官 RBC haemolysis 和 EC death \\
Septic shock & 多器官 RBC haemolysis 和 EC death \\
\bottomrule
\end{longtable}

\subsubsection{小鼠模型的發現}

\textbf{COVID-19 小鼠模型 vs 缺血模型}：
\begin{itemize}
    \item SARS-CoV-2 感染健康年輕小鼠 5 天：主要在心臟微血管有 RBC haemolysis
    \item 肺和腦有明顯損傷，但無 RBC haemolysis
    \item 結論：病毒感染本身和相關器官損傷不足以誘導 RBC haemolysis
\end{itemize}

\textbf{心臟缺血模型}：
\begin{itemize}
    \item LAD ligation（ischaemia-reperfusion, IR）：廣泛 RBC haemolysis
    \item Permanent ligation (PL)：同樣有 RBC haemolysis
    \item 結論：\textbf{Ischaemia 本身足以誘導 RBC haemolysis}
\end{itemize}

\subsection{兩種不同的 EC 死亡途徑}

研究發現兩種空間上不同的 EC 死亡模式，各自調控不同的微血管阻塞機制：

\begin{longtable}{p{3.5cm}p{5.5cm}p{5cm}}
\toprule
\textbf{特徵} & \textbf{Necroptosis (PS$^{low}$)} & \textbf{Apoptosis (PS$^{high}$)} \\
\midrule
\endhead
位置 & Villus 尖端（apical region） & Villus 下部（lower body） \\
Annexin V 結合 & 弱（PS$^{low}$） & 強（PS$^{high}$） \\
PI/RedDot2 通透性 & 陽性（細胞膜破壞） & 陰性 \\
TUNEL & 陰性 & 陽性 \\
Hoechst 染色 & 減少（核溶解） & 增強/濃縮 \\
Caspase-3 & 陰性 & 陽性（cleaved caspase-3） \\
p-MLKL/p-RIPK3 & 陽性 & 陰性 \\
導致的阻塞 & \textbf{RBC haemolysis} & \textbf{Platelet + Fibrin} \\
佔 PS$^+$ EC 比例 & 約 70\% & 約 30\% \\
\bottomrule
\end{longtable}

\begin{importantbox}
\textbf{關鍵概念}：
\begin{itemize}
    \item \textbf{Necroptotic EC} $\rightarrow$ RBC haemolysis $\rightarrow$ RBC-dependent 微血管阻塞（佔 $>$ 65\%）
    \item \textbf{Apoptotic EC} $\rightarrow$ Platelet + Fibrin 沉積 $\rightarrow$ 傳統微血栓
\end{itemize}
這兩種機制在空間上不重疊！
\end{importantbox}

\subsection{EC Necroptosis 和 Complement 介導 Haemolysis}

\subsubsection{MLKL 的角色}

MLKL (Mixed Lineage Kinase Domain-Like pseudokinase) 是 necroptosis 途徑的終端媒介：

\begin{itemize}
    \item \textbf{Mlkl$^{-/-}$ 小鼠}：RBC haemolysis 明顯減少
    \item \textbf{Mlkl$^{EC-/-}$ 小鼠}：RBC haemolysis 減少（但比全身 KO 少）
    \item \textbf{Mlkl$^{RBC-/-}$ 小鼠}：RBC haemolysis 正常
\end{itemize}

\note{這證明 \textbf{EC necroptosis}（而非 RBC necroptosis）是促進 RBC haemolysis 的主要機制。}

\subsubsection{Complement 的角色}

垂死的 EC 促進 complement cascade 活化：
\begin{itemize}
    \item C9$^{-/-}$ 小鼠（無法形成 membrane attack complex, MAC, C5b-9）：RBC haemolysis 明顯減少
    \item IF 確認 C5b-9 沉積在溶血 RBC 表面，但不在完整 RBC 上
    \item 人類 COVID-19 和缺血性器官中也觀察到類似的 C9 沉積
\end{itemize}

\subsection{溶血 RBC 形成止血栓塞（Endo-seal）}

\subsubsection{Endo-seal 的形成過程}

Intravital microscopy 觀察到的動態過程：
\begin{enumerate}
    \item RBC 黏附到垂死 EC 表面（PS$^+$）
    \item 黏附的 RBC 在血流拖曳力下變形伸長
    \item RBC membrane 碎裂
    \item Membrane fragments 沉積在血管壁
    \item 沉積的 fragments 作為後續 RBC 的 docking sites
    \item 重複循環 $\rightarrow$ 多層 RBC membrane 堆積
    \item 形成連續的 endovascular sealant（``endo-seal''，厚度 5--10 $\mu$m）
\end{enumerate}

\subsubsection{Endo-seal 的止血功能}

\begin{itemize}
    \item RBC haemolysis 程度與 interstitial bleeding 呈\textbf{負相關}（R = -0.6577, p $<$ 0.0001）
    \item 有 endo-seal 的區域：無出血
    \item 缺乏 endo-seal 的區域：可見 RBC extravasation（出血）
\end{itemize}

\textbf{基因剔除證據}：
\begin{itemize}
    \item Mlkl$^{-/-}$、Mlkl$^{EC-/-}$、C9$^{-/-}$ 小鼠：haemolysis 減少 $\rightarrow$ 微血管出血增加
    \item 此止血機制不受 platelet depletion、potent platelet inhibitors 或 high-dose anticoagulation 影響
\end{itemize}

\begin{importantbox}
\textbf{重大發現}：Endo-seal 是一種全新的止血機制：
\begin{itemize}
    \item 獨立於 platelet 和 fibrin
    \item 可在低流速（$<$ 25 s$^{-1}$）和酸性環境（pH 6.0--7.0）下運作
    \item 這是傳統 platelet 和 coagulation 反應受抑制的環境
\end{itemize}
\end{importantbox}

\subsection{溶血 RBC 促進 RBC 聚集和阻塞}

\subsubsection{Homotypic RBC Aggregation}

溶血 RBC membranes 具有黏附特性：
\begin{itemize}
    \item 在完整 RBC 之間形成 adhesive bridges
    \item RBC aggregation 發生在極低 shear rate（$<$ 25 s$^{-1}$）
    \item 不需要 platelet、neutrophil 或 plasma proteins
    \item 這是 RBC membrane 的內在特性
\end{itemize}

\subsubsection{In Vitro 驗證}

\begin{itemize}
    \item 將全血灌注到 necrotic HMEC-1（human microvascular EC）表面
    \item 足以誘導 RBC aggregation
    \item Aggregates 通過受損 RBC membranes 錨定到 HMEC-1
    \item C9 沉積在 necrotic HMEC-1 上
    \item 抑制 complement 活化（Cp40）可防止 RBC aggregation
\end{itemize}

\subsubsection{微血管阻塞的兩種機制}

\begin{longtable}{p{4cm}p{5cm}p{5cm}}
\toprule
\textbf{參數} & \textbf{RBC Microangiopathy} & \textbf{Platelet-Fibrin Microthrombi} \\
\midrule
\endhead
Shear rate & $<$ 25 s$^{-1}$ & $>$ 50 s$^{-1}$ \\
觸發條件 & Ischaemia, acidosis & Reperfusion \\
EC 死亡類型 & Necroptosis & Apoptosis \\
主要成分 & 溶血 RBC membranes & Platelet + Fibrin \\
需要 complement & 是（C5b-9/MAC） & 否 \\
位置 & Villus tip & Villus base \\
\bottomrule
\end{longtable}

\subsection{Mlkl$^{EC-/-}$ 小鼠的器官保護}

\begin{itemize}
    \item EC-specific Mlkl deletion 顯著減少 gut IR tissue injury
    \item 減少溶血 RBC 促進的 RBC aggregates
    \item 改善 FITC-dextran 灌注
    \item Laser speckle contrast imaging 顯示組織灌流改善
\end{itemize}

\note{這表明 \textbf{EC necroptosis 和相關的 RBC haemolysis 是微血管阻塞和局部組織損傷的關鍵因素}，為治療提供了新標靶。}

%============================================================
\section{機制總結}
%============================================================

\subsection{RBC Microangiopathy 的病理生理機制}

\begin{enumerate}
    \item \textbf{Ischaemia} $\rightarrow$ 組織缺氧 + 代謝性酸中毒（pH 6.0--7.0）
    \item \textbf{EC injury} $\rightarrow$ MLKL phosphorylation $\rightarrow$ \textbf{EC necroptosis}
    \item \textbf{Dying EC} $\rightarrow$ 暴露 phosphatidylserine (PS) $\rightarrow$ Complement 活化
    \item \textbf{C5b-9 (MAC) formation} $\rightarrow$ RBC membrane damage $\rightarrow$ \textbf{Haemolysis}
    \item \textbf{Haemolysed RBC membranes}:
    \begin{itemize}
        \item 沉積在 EC 表面 $\rightarrow$ 形成 \textbf{endo-seal}（止血）
        \item 在 RBC 之間形成 adhesive bridges $\rightarrow$ \textbf{RBC aggregation}（阻塞）
    \end{itemize}
    \item \textbf{過度反應} $\rightarrow$ 微血管阻塞 $\rightarrow$ 組織缺血加劇
\end{enumerate}

\subsection{為何傳統抗凝血治療效果有限？}

\begin{itemize}
    \item 傳統治療針對 platelet 和 fibrin
    \item 但 COVID-19 微血管阻塞主要由 \textbf{RBC haemolysis} 造成
    \item RBC microangiopathy 在低 shear rate 和 acidosis 條件下發生
    \item 這些條件下 platelet function 和 coagulation 反應本就受抑制
    \item 因此抗 platelet 和 anticoagulant 無法有效改善此類阻塞
\end{itemize}

%============================================================
\section{臨床意義}
%============================================================

\subsection{疾病相關性}

此 RBC haemostatic/microangiopathy 機制與多種疾病相關：

\begin{longtable}{p{5cm}p{9cm}}
\toprule
\textbf{疾病} & \textbf{相關性} \\
\midrule
\endhead
COVID-19（急性） & 多器官微血管 RBC haemolysis 和阻塞 \\
Long COVID & 持續的 capillary dysfunction、循環 haem 升高與 terminal complement complex 形成相關 \\
Acute myocardial infarction & 缺血區域的 RBC haemolysis 和 no-reflow phenomenon \\
Ischaemic stroke & 腦微血管 RBC haemolysis \\
Septic shock & 多器官微血管 RBC microangiopathy \\
Cardiogenic shock & 多器官微血管 RBC microangiopathy \\
Acute coronary syndromes & Microvascular dysfunction \\
\bottomrule
\end{longtable}

\subsection{潛在治療策略}

\begin{itemize}
    \item \textbf{Necroptosis 抑制劑}：減少 EC necroptosis
    \item \textbf{Complement 抑制}：防止 RBC haemolysis
    \item \textbf{Haemopexin}：中和過量游離 haem
    \item \textbf{聯合治療}：可能更有效
\end{itemize}

\note{
\textbf{治療考量}：
\begin{itemize}
    \item 抑制 necroptosis 可能增加微血管出血風險（因 endo-seal 減少）
    \item 類似於抗 platelet/anticoagulant 的出血風險
    \item 需要仔細評估 benefit-risk ratio
\end{itemize}
}

%============================================================
\section{研究局限與未來方向}
%============================================================

\subsection{研究局限}

\begin{itemize}
    \item 主要基於屍檢組織和動物模型
    \item 人體內的動態過程難以直接觀察
    \item 治療效果需要臨床試驗驗證
\end{itemize}

\subsection{未來研究方向}

\begin{itemize}
    \item 開發針對 RBC microangiopathy 的特異性治療
    \item 評估聯合抑制 necroptosis + complement 的效果和安全性
    \item 研究 long COVID 中此機制的持續性
    \item 探索其他缺血再灌注疾病中的應用
\end{itemize}

%============================================================
\section{名詞解釋}
%============================================================

\begin{longtable}{p{5cm}p{9cm}}
\toprule
\textbf{名詞} & \textbf{解釋} \\
\midrule
\endhead
Necroptosis & 一種可調控的細胞壞死形式，由 RIPK1-RIPK3-MLKL 途徑介導 \\
Apoptosis & 程序性細胞凋亡，由 caspase 介導，不引起發炎 \\
MLKL & Mixed Lineage Kinase Domain-Like pseudokinase，necroptosis 的終端執行者 \\
RIPK1/RIPK3 & Receptor-Interacting Protein Kinase 1/3，necroptosis 途徑的上游激酶 \\
Phosphatidylserine (PS) & 正常位於細胞膜內層的磷脂，細胞死亡時外翻暴露 \\
Annexin V & 結合 PS 的蛋白，用於檢測細胞死亡 \\
C5b-9 (MAC) & Membrane Attack Complex，complement 終末產物，在細胞膜上形成孔洞 \\
Glycophorin A (CD235) & RBC 表面的主要唾液酸糖蛋白 \\
VE-cadherin & Vascular endothelial cadherin，EC 間的黏附分子 \\
Endo-seal & 由溶血 RBC membrane 形成的內血管止血栓塞 \\
\bottomrule
\end{longtable}

%============================================================
\section{重點摘要}
%============================================================

\begin{importantbox}
\textbf{Take-Home Messages}：
\begin{enumerate}
    \item COVID-19 微血管阻塞的主要原因是 \textbf{RBC haemolysis}，而非傳統認為的 platelet-fibrin microthrombi
    \item 機制：\textbf{Ischaemia} $\rightarrow$ \textbf{EC necroptosis (MLKL)} $\rightarrow$ \textbf{Complement activation (C5b-9)} $\rightarrow$ \textbf{RBC haemolysis}
    \item 溶血 RBC membrane 具有雙重功能：
    \begin{itemize}
        \item \textbf{Endo-seal}：防止微血管出血（有益）
        \item \textbf{RBC aggregation}：造成微血管阻塞（有害）
    \end{itemize}
    \item 這是一種\textbf{全新的止血機制}，獨立於 platelet 和 fibrin，可在低流速和酸性環境下運作
    \item 此機制不僅見於 COVID-19，也見於 \textbf{MI、stroke、sepsis、cardiogenic shock} 等缺血再灌注相關疾病
    \item 治療標靶：\textbf{Necroptosis 抑制} + \textbf{Complement 抑制} + \textbf{Haemopexin}
    \item 這解釋了為何傳統 anticoagulant/antiplatelet 治療對 COVID-19 微血管功能障礙效果有限
\end{enumerate}
\end{importantbox}

%============================================================
\section{參考文獻}
%============================================================

\begin{enumerate}
    \item Wu MCL, Italiano E, Jarvis-Child R, et al. Ischaemic endothelial necroptosis induces haemolysis and COVID-19 angiopathy. \href{https://doi.org/10.1038/s41586-025-09076-x}{\textit{Nature}. 2025. DOI: 10.1038/s41586-025-09076-x}

    \item Flaumenhaft R, Enjyoji K, Schmaier AA. Vasculopathy in COVID-19. \href{https://doi.org/10.1182/blood.2021014181}{\textit{Blood}. 2022;140:222-235.}

    \item Conway EM, et al. Understanding COVID-19-associated coagulopathy. \href{https://doi.org/10.1038/s41577-022-00762-9}{\textit{Nat Rev Immunol}. 2022;22:639-649.}

    \item Ackermann M, et al. Pulmonary vascular endothelialitis, thrombosis, and angiogenesis in COVID-19. \href{https://doi.org/10.1056/NEJMoa2015432}{\textit{N Engl J Med}. 2020;383:120-128.}

    \item Pasparakis M, Vandenabeele P. Necroptosis and its role in inflammation. \href{https://doi.org/10.1038/nature14191}{\textit{Nature}. 2015;517:311-320.}

    \item Afzali B, Noris M, Lambrecht BN, Kemper C. The state of complement in COVID-19. \href{https://doi.org/10.1038/s41577-021-00665-1}{\textit{Nat Rev Immunol}. 2022;22:77-84.}

    \item Cervia-Hasler C, et al. Persistent complement dysregulation with signs of thromboinflammation in active long Covid. \href{https://doi.org/10.1126/science.adg7942}{\textit{Science}. 2024;383:eadg7942.}
\end{enumerate}

\end{document}
